
\textbf{Setup.} The local parameters are
\(\hat{\bm{e}} = (\epsilon_0, \bm{\beta}, \tau, \bm{\epsilon}_1)\) with
sizes \((1, 2, 1, 2)\). The 3D strain is:
\[\bm{\varepsilon} = \varepsilon\,\FrameBasis \otimes \FrameBasis + \FrameBasis \stackrel{s}{\otimes} \bm{\varepsilon}_\gamma - \nu\varepsilon\,\mathbf{1}_\Omega\]
where:
\[\varepsilon = \epsilon_0 + \bm{r} \cdot \bm{\epsilon}_1, \qquad \bm{\varepsilon}_\gamma = \tau\,\bm{\tau}_\varphi + \bm{\tau}_\psi\,\bm{\beta}\]
with:
\[\bm{\tau}_\varphi := \FrameBasis \times \bm{r} + \nabla\WarpTwistIesan, \qquad \bm{\tau}_\psi := \nabla\bm{\WarpShearIesan} + \bm{\Pi}_{(b)}\]

\textbf{Computing
\(\bar{\mathbf{C}} = \int_\Section \mathbf{B}^T \mathbb{C} \mathbf{B} \dA\).}

The strain energy density splits into axial and shear parts:
\[\tfrac{1}{2}\bm{\varepsilon}^T \mathbb{C} \bm{\varepsilon} = \tfrac{1}{2}E\varepsilon^2 + \tfrac{1}{2}G|\bm{\varepsilon}_\gamma|^2\]

In the partition \((\epsilon_0, \bm{\beta}, \tau, \bm{\epsilon}_1)\):

\[\bar{\mathbf{C}} = \begin{pmatrix}
EA & \mathbf{0}^\top & 0 & EA\CentroidLocation^\top \\[4pt]
\mathbf{0} & G\mathbf{H}_{\psi\psi} & G\mathbf{h}_{\varphi\psi} & \mathbf{0} \\[4pt]
0 & G\mathbf{h}_{\varphi\psi}^\top & G\Jsv & \mathbf{0}^\top \\[4pt]
EA\CentroidLocation & \mathbf{0} & \mathbf{0} & E\IntRoR
\end{pmatrix}\]

where:
\[\Jsv := \int_\Section |\bm{\tau}_\varphi|^2 \dA, \qquad \mathbf{h}_{\varphi\psi} := \int_\Section \bm{\tau}_\psi^\top \bm{\tau}_\varphi \dA, \qquad \mathbf{H}_{\psi\psi} := \int_\Section \bm{\tau}_\psi^\top \bm{\tau}_\psi \dA\]

\textbf{Evaluating \(\mathbf{h}_{\varphi\psi}\).}

We have:
\[\mathbf{h}_{\varphi\psi} = \int_\Section (\nabla\bm{\WarpShearIesan} + \bm{\Pi}_{(b)})^\top (\FrameBasis \times \bm{r} + \nabla\WarpTwistIesan) \dA\]

Breaking into four terms:

\emph{Term 1:}
\(\int_\Section (\nabla\bm{\WarpShearIesan})^\top (\FrameBasis \times \bm{r}) \dA\)

Using divergence theorem with
\(\nabla \cdot (\FrameBasis \times \bm{r}) = 0\):
\[\int_\Section (\FrameBasis \times \bm{r}) \cdot \nabla\WarpShearIesan_\alpha \dA = \oint_{\partial\Section} \WarpShearIesan_\alpha (\FrameBasis \times \bm{r}) \cdot \bm{\nu} \, ds\]

\emph{Term 2:}
\(\int_\Section (\nabla\bm{\WarpShearIesan})^\top \nabla\WarpTwistIesan \dA\)

Using Green's identity with \(\Delta\WarpTwistIesan = 0\) and
\(\partial_\nu\WarpTwistIesan = -(\FrameBasis \times \bm{r}) \cdot \bm{\nu}\):
\[\int_\Section \nabla\WarpTwistIesan \cdot \nabla\WarpShearIesan_\alpha \dA = \oint_{\partial\Section} \WarpShearIesan_\alpha \partial_\nu\WarpTwistIesan \, ds = -\oint_{\partial\Section} \WarpShearIesan_\alpha (\FrameBasis \times \bm{r}) \cdot \bm{\nu} \, ds\]

Terms 1 and 2 cancel:
\[\int_\Section (\nabla\bm{\WarpShearIesan})^\top (\FrameBasis \times \bm{r} + \nabla\WarpTwistIesan) \dA = \mathbf{0}\]

\emph{Term 3:}
\(\int_\Section \bm{\Pi}_{(b)}^\top (\FrameBasis \times \bm{r}) \dA\)

With
\(\bm{\Pi}_{(b)} = -\nu(\text{dev}(\bm{r} \otimes \bm{r}) - \bm{r} \otimes \CentroidLocation)\),
acting on column \(\alpha\):
\[(\bm{\Pi}_{(b)})_\alpha = -\nu\big((r_\alpha - \bar{\rho}_\alpha)\bm{r} - \tfrac{1}{2}|\bm{r}|^2 \bm{e}_\alpha\big)\]

Then:
\[(\FrameBasis \times \bm{r}) \cdot (\bm{\Pi}_{(b)})_\alpha = -\nu(r_\alpha - \bar{\rho}_\alpha)\underbrace{(\FrameBasis \times \bm{r}) \cdot \bm{r}}_{=0} + \tfrac{\nu}{2}|\bm{r}|^2 (\FrameBasis \times \bm{r})_\alpha\]

So:
\[\int_\Section \bm{\Pi}_{(b)}^\top (\FrameBasis \times \bm{r}) \dA = \tfrac{\nu}{2} \int_\Section |\bm{r}|^2 (\FrameBasis \times \bm{r}) \dA\]

\emph{Term 4:}
\(\int_\Section \bm{\Pi}_{(b)}^\top \nabla\WarpTwistIesan \dA\)

\[\int_\Section (\bm{\Pi}_{(b)})_\alpha \cdot \nabla\WarpTwistIesan \dA = -\nu \int_\Section (r_\alpha - \bar{\rho}_\alpha)(\bm{r} \cdot \nabla\WarpTwistIesan) \dA + \tfrac{\nu}{2} \int_\Section |\bm{r}|^2 \partial_\alpha\WarpTwistIesan \dA\]

Combining terms 3 and 4:
\[\mathbf{h}_{\varphi\psi} = \tfrac{\nu}{2} \int_\Section |\bm{r}|^2 (\FrameBasis \times \bm{r} + \nabla\WarpTwistIesan) \dA - \nu \int_\Section (\bm{r} - \CentroidLocation)(\bm{r} \cdot \nabla\WarpTwistIesan) \dA\]

Define:
\[\bm{\mu}_\varphi := \int_\Section |\bm{r}|^2 \bm{\tau}_\varphi \dA, \qquad \bm{\lambda}_\varphi := \int_\Section (\bm{r} \cdot \nabla\WarpTwistIesan)(\bm{r} - \CentroidLocation) \dA\]

Then:
\[\mathbf{h}_{\varphi\psi} = \tfrac{\nu}{2}\bm{\mu}_\varphi - \nu\bm{\lambda}_\varphi\]

\textbf{Relation to shear center.} The shear center is defined by:
\[\Jsv\,\CenterShearLocation = -\int_\Section \bm{r} \times \bm{\tau}_\psi \cdot \FrameBasis \dA = \int_\Section (\FrameBasis \times \bm{r})^\top \bm{\tau}_\psi \dA\]

This is \emph{not} the same as \(\mathbf{h}_{\varphi\psi}\) because the
latter includes the \(\nabla\WarpTwistIesan\) term.

\subsubsection{The energy-conjugacy
condition.}\label{the-energy-conjugacy-condition.}

For beam strains \(\bm{e} = \mathbf{X}\hat{\bm{e}}\) to be
work-conjugate to resultants
\(\bm{s} = \mathbf{R}\mathbf{L}^{-1}\hat{\bm{e}}\):
\[\bm{s}^\top \delta\bm{e} = \int_\Section \bm{\sigma}^\top \delta\bm{\varepsilon} \dA = \hat{\bm{e}}^\top \bar{\mathbf{C}} \delta\hat{\bm{e}}\]

This requires:
\[\mathbf{X}^\top \mathbf{R} \mathbf{L}^{-1} = \bar{\mathbf{C}}\]

Hence:
\[\mathbf{X}^{-1} = \bar{\mathbf{C}}^{-1} \mathbf{L}^{-\top} \mathbf{R}^\top\]

\textbf{Computing \(\mathbf{L}^{-\top} \mathbf{R}^\top\).}

Recall (with partition
\(\hat{\bm{e}} = (\epsilon_0, \bm{\beta}, \tau, \bm{\epsilon}_1)\) and
\(\bm{p} = (a_\varepsilon, \bm{a}_\Omega, a_\vartheta, \bm{b}_{\Section})\)):

\[\mathbf{L}_0 = \begin{pmatrix}
1 & \mathbf{0}^\top & 0 & \mathbf{0}^\top \\
\mathbf{0} & \mathbf{0} & \mathbf{0} & \oneS \\
0 & \mathbf{0}^\top & 1 & \CenterShearLocation^\top \\
\mathbf{0} & \oneS & \mathbf{0} & \mathbf{0}
\end{pmatrix}
\qquad\text{so}\qquad 
\mathbf{L}_0^{-\top} = \begin{pmatrix}
1 & \mathbf{0}^\top & 0 & \mathbf{0}^\top \\
\mathbf{0} & \mathbf{0} & -\CenterShearLocation & \oneS \\
0 & \mathbf{0}^\top & 1 & \mathbf{0}^\top \\
\mathbf{0} & \oneS & \mathbf{0} & \mathbf{0}
\end{pmatrix}\]

And with partition \(\bm{s} = (N, \bm{F}, T, \bm{M})\) (but noting
\(\bm{M}\) appears as \(\FrameBasis \times \bm{M}\) in the formulas):

\[\mathbf{R}^\top = \begin{pmatrix}
-EA & \mathbf{0}^\top & 0 & -EA(\FrameBasis \times \CentroidLocation)^\top \\
-EA\CentroidLocation & \mathbf{0} & \mathbf{0} & -(\FrameBasis \times E\IntRoR)^\top \\
0 & \mathbf{0}^\top & -G\Jsv & \mathbf{0}^\top \\
\mathbf{0} & -E\IntRGoRG & \mathbf{0} & \mathbf{0}
\end{pmatrix}\]


\paragraph{Energy matrix \(\bar{\mathbf{C}}\).}

In partition
\(\hat{\bm{e}} = (\epsilon_0, \bm{\beta}, \tau, \bm{\epsilon}_1)\):
\[\bar{\mathbf{C}} = \begin{pmatrix}
EA & \mathbf{0}^\top & 0 & EA\CentroidLocation^\top \\
\mathbf{0} & G\mathbf{H}_{\psi\psi} & G\mathbf{h}_{\varphi\psi} & \mathbf{0} \\
0 & G\mathbf{h}_{\varphi\psi}^\top & G\Jsv & \mathbf{0}^\top \\
EA\CentroidLocation & \mathbf{0} & \mathbf{0} & E\IntRoR
\end{pmatrix}
\]

\begin{Details}
Verifying the \((3,3)\) entry (torsional stiffness):

\[
\int_\Section |\bm{\tau}_\varphi|^2 \dA = \int_\Section |\FrameBasis \times \bm{r} + \nabla\WarpTwistIesan|^2 \dA
\]
\[= \int_\Section r^2 \dA + 2\int_\Section (\FrameBasis \times \bm{r}) \cdot \nabla\WarpTwistIesan \dA + \int_\Section |\nabla\WarpTwistIesan|^2 \dA
\]

For the last term, using Green's first identity with
\(\Delta\WarpTwistIesan = 0\):
\[\int_\Section |\nabla\WarpTwistIesan|^2 \dA = \oint_{\partial\Section} \WarpTwistIesan \,\partial_\nu\WarpTwistIesan \, ds = -\oint_{\partial\Section} \WarpTwistIesan (\FrameBasis \times \bm{r}) \cdot \bm{\nu} \, ds\]

By divergence theorem (using
\(\nabla \cdot (\FrameBasis \times \bm{r}) = 0\)):
\[\oint_{\partial\Section} \WarpTwistIesan (\FrameBasis \times \bm{r}) \cdot \bm{\nu} \, ds = \int_\Section \nabla\WarpTwistIesan \cdot (\FrameBasis \times \bm{r}) \dA\]

Therefore:
\[\int_\Section |\nabla\WarpTwistIesan|^2 \dA = -\int_\Section (\FrameBasis \times \bm{r}) \cdot \nabla\WarpTwistIesan \dA\]

So:
\[\int_\Section |\bm{\tau}_\varphi|^2 \dA = \int_\Section r^2 \dA + \int_\Section (\FrameBasis \times \bm{r}) \cdot \nabla\WarpTwistIesan \dA = \Jsv\]

\textbf{The (2,3) entry (shear-torsion coupling):}

\[\mathbf{h}_{\varphi\psi} = \int_\Section \bm{\tau}_\psi^\top \bm{\tau}_\varphi \dA = \int_\Section (\bm{\Pi}_{(b)}^\top + \nabla\bm{\WarpShearIesan})(\FrameBasis \times \bm{r} + \nabla\WarpTwistIesan) \dA\]

Breaking into four terms:

\emph{Term A:}
\(\int_\Section (\nabla\bm{\WarpShearIesan})(\FrameBasis \times \bm{r}) \dA\)

\emph{Term B:}
\(\int_\Section (\nabla\bm{\WarpShearIesan})\nabla\WarpTwistIesan \dA\)

Using Green's identity with \(\Delta\WarpTwistIesan = 0\) and
\(\partial_\nu\WarpTwistIesan = -(\FrameBasis \times \bm{r}) \cdot \bm{\nu}\):
\[[\text{Term B}]_\alpha = \oint_{\partial\Section} \WarpShearIesan_\alpha \,\partial_\nu\WarpTwistIesan \, ds = -\oint_{\partial\Section} \WarpShearIesan_\alpha (\FrameBasis \times \bm{r}) \cdot \bm{\nu} \, ds\]

By divergence theorem:
\[[\text{Term A}]_\alpha = \oint_{\partial\Section} \WarpShearIesan_\alpha (\FrameBasis \times \bm{r}) \cdot \bm{\nu} \, ds\]

So Terms A and B cancel:
\(\int_\Section (\nabla\bm{\WarpShearIesan})\bm{\tau}_\varphi \dA = \mathbf{0}\).
✓

\emph{Term C:}
\(\int_\Section \bm{\Pi}_{(b)}^\top (\FrameBasis \times \bm{r}) \dA\)

\emph{Term D:}
\(\int_\Section \bm{\Pi}_{(b)}^\top \nabla\WarpTwistIesan \dA\)

For the \(\alpha\)-th component, with
\((\bm{\Pi}_{(b)})_\alpha = -\nu((r_\alpha - \bar{\rho}_\alpha)\bm{r} - \frac{1}{2}r^2\bm{e}_\alpha)\):

\textbf{Term C:}
\[(\bm{\Pi}_{(b)})_\alpha \cdot (\FrameBasis \times \bm{r}) = -\nu(r_\alpha - \bar{\rho}_\alpha)\underbrace{\bm{r} \cdot (\FrameBasis \times \bm{r})}_{=0} + \frac{\nu}{2}r^2(\FrameBasis \times \bm{r})_\alpha\]

So:
\[\text{Term C} = \frac{\nu}{2}\int_\Section r^2 (\FrameBasis \times \bm{r}) \dA\]

\textbf{Term D:}
\[(\bm{\Pi}_{(b)})_\alpha \cdot \nabla\WarpTwistIesan = -\nu(r_\alpha - \bar{\rho}_\alpha)(\bm{r} \cdot \nabla\WarpTwistIesan) + \frac{\nu}{2}r^2 \partial_\alpha\WarpTwistIesan\]

So:
\[\text{Term D} = \frac{\nu}{2}\int_\Section r^2 \nabla\WarpTwistIesan \dA - \nu\int_\Section (\bm{r} - \CentroidLocation)(\bm{r} \cdot \nabla\WarpTwistIesan) \dA\]

Therefore:
\[\mathbf{h}_{\varphi\psi} = \frac{\nu}{2}\int_\Section r^2 \bm{\tau}_\varphi \dA - \nu\int_\Section (\bm{r} - \CentroidLocation)(\bm{r} \cdot \nabla\WarpTwistIesan) \dA\]

\textbf{Relation to shear center.} The shear center satisfies:
\[\Jsv\,\CenterShearLocation = \int_\Section \bm{\tau}_\psi^\top (\bm{r} \times \FrameBasis) \dA = -\int_\Section \bm{\tau}_\psi^\top (\FrameBasis \times \bm{r}) \dA\]

This involves only \((\FrameBasis \times \bm{r})\), not the full
\(\bm{\tau}_\varphi = \FrameBasis \times \bm{r} + \nabla\WarpTwistIesan\).
So: \[\mathbf{h}_{\varphi\psi} \neq -\Jsv\,\CenterShearLocation\]

The difference is precisely the \(\nabla\WarpTwistIesan\) coupling
terms. Define:
\[\bm{\lambda} := \int_\Section (\bm{r} - \CentroidLocation)(\bm{r} \cdot \nabla\WarpTwistIesan) \dA, \qquad \bm{\zeta} := \int_\Section r^2 \nabla\WarpTwistIesan \dA\]

Then:
\[\mathbf{h}_{\varphi\psi} = -\Jsv\,\CenterShearLocation + \frac{\nu}{2}\bm{\zeta} - \nu\bm{\lambda}\]
\end{Details}

The quantities \(\bm{\zeta}\) and \(\bm{\lambda}\) encode torsion-shear
coupling that exists even for symmetric sections when the origin is not
at the centroid. These do NOT vanish under the normalization
\(\int \WarpTwistIesan \dA = 0\).

\subsubsection{Derive X}

\textbf{Step 1: Compute \(\mathbf{L}_0^{-\top}\).}

We have: \[\mathbf{L}_0^{-1} = \begin{pmatrix}
1 & \mathbf{0}^\top & 0 & \mathbf{0}^\top \\
\mathbf{0} & \mathbf{0} & \mathbf{0} & \oneS \\
0 & -\CenterShearLocation^\top & 1 & \mathbf{0}^\top \\
\mathbf{0} & \oneS & \mathbf{0} & \mathbf{0}
\end{pmatrix}
\qquad
\mathbf{L}_0^{-\top} = \begin{pmatrix}
1 & \mathbf{0}^\top & 0 & \mathbf{0}^\top \\
\mathbf{0} & \mathbf{0} & -\CenterShearLocation & \oneS \\
0 & \mathbf{0}^\top & 1 & \mathbf{0}^\top \\
\mathbf{0} & \oneS & \mathbf{0} & \mathbf{0}
\end{pmatrix}\]

\textbf{Step 2: Compute \(\mathbf{R}^\top\).}

From the resultant matrix with partition
\(\bm{s} = (N, \bm{F}, T, \bm{M})\),
\(\bm{p} = (a_\varepsilon, \bm{a}_\Omega, a_\vartheta, \bm{b}_{\Section})\):

\[\mathbf{R}^\top = \begin{pmatrix}
EA & \mathbf{0}^\top & 0 & EA(\FrameBasis\times\CentroidLocation)^\top \\
EA\CentroidLocation & \mathbf{0} & \mathbf{0} & E(\FrameBasis^\times\IntRoR)^\top \\
0 & \mathbf{0}^\top & G\Jsv & \mathbf{0}^\top \\
\mathbf{0} & E\IntRGoRG & \mathbf{0} & \mathbf{0}
\end{pmatrix}\]

Using
\((\FrameBasis\times\CentroidLocation)^\top = -\CentroidLocation^\top\FrameBasis^\times\)
and \((\FrameBasis^\times\IntRoR)^\top = -\IntRoR\FrameBasis^\times\):

\[\mathbf{R}^\top = \begin{pmatrix}
EA & \mathbf{0}^\top & 0 & -EA\CentroidLocation^\top\FrameBasis^\times \\
EA\CentroidLocation & \mathbf{0} & \mathbf{0} & -E\IntRoR\FrameBasis^\times \\
0 & \mathbf{0}^\top & G\Jsv & \mathbf{0}^\top \\
\mathbf{0} & E\IntRGoRG & \mathbf{0} & \mathbf{0}
\end{pmatrix}
\]

\subparagraph{Step 3: Compute \(\mathbf{L}_0^{-\top}\mathbf{R}^\top\).}

\[
\mathbf{L}_0^{-\top}\mathbf{R}^\top = \begin{pmatrix}
EA & \mathbf{0}^\top & 0 & -EA\CentroidLocation^\top\FrameBasis^\times \\
\mathbf{0} & E\IntRGoRG & -G\Jsv\CenterShearLocation & \mathbf{0} \\
0 & \mathbf{0}^\top & G\Jsv & \mathbf{0}^\top \\
EA\CentroidLocation & \mathbf{0} & \mathbf{0} & -E\IntRoR\FrameBasis^\times
\end{pmatrix}
\]

\subparagraph{Step 4: Invert \(\bar{\mathbf{C}}\).}

The energy matrix has block structure:
\[\bar{\mathbf{C}} = \begin{pmatrix}
EA & \mathbf{0}^\top & 0 & EA\CentroidLocation^\top \\
\mathbf{0} & G\mathbf{H}_{\psi\psi} & G\mathbf{h}_{\varphi\psi} & \mathbf{0} \\
0 & G\mathbf{h}_{\varphi\psi}^\top & G\Jsv & \mathbf{0}^\top \\
EA\CentroidLocation & \mathbf{0} & \mathbf{0} & E\IntRoR
\end{pmatrix}\]

Reordering to \((\epsilon_0, \bm{\epsilon}_1, \bm{\beta}, \tau)\)
reveals block-diagonal structure:

\[
\bar{\mathbf{C}}_{\rm perm} = \begin{pmatrix}
\mathbf{A} & \mathbf{0} \\
\mathbf{0} & \mathbf{D}
\end{pmatrix}
\qquad\text{with:}\qquad
\mathbf{A} = \begin{pmatrix} EA & EA\CentroidLocation^\top \\ EA\CentroidLocation & E\IntRoR 
\end{pmatrix}, 
\quad 
\mathbf{D} = G\begin{pmatrix} \mathbf{H}_{\psi\psi} & \mathbf{h}_{\varphi\psi} \\ \mathbf{h}_{\varphi\psi}^\top & \Jsv \end{pmatrix}
\]

\textbf{Inverting \(\mathbf{A}\):} Schur complement is
\(E\IntRoR - EA\CentroidLocation\otimes\CentroidLocation = E\IntRGoRG\):
\[\mathbf{A}^{-1} = \begin{pmatrix} 
\frac{1}{EA} + \CentroidLocation^\top(E\IntRGoRG)^{-1}\CentroidLocation & -\CentroidLocation^\top(E\IntRGoRG)^{-1} \\
-(E\IntRGoRG)^{-1}\CentroidLocation & (E\IntRGoRG)^{-1}
\end{pmatrix},
\qquad\text{and}\qquad 
\mathbf{D}^{-1} = \frac{1}{G}\begin{pmatrix}
\mathbf{H}_{\psi\psi}^{-1} + \frac{1}{\tilde{\Jsv}}\bm{\eta}\otimes\bm{\eta} & -\frac{1}{\tilde{\Jsv}}\bm{\eta} \\
-\frac{1}{\tilde{\Jsv}}\bm{\eta}^\top & \frac{1}{\tilde{\Jsv}}
\end{pmatrix}\]

\[\tilde{\Jsv} := \Jsv - \mathbf{h}_{\varphi\psi}^\top\mathbf{H}_{\psi\psi}^{-1}\mathbf{h}_{\varphi\psi}
\qquad\text{and}\qquad 
\bm{\eta} := \mathbf{H}_{\psi\psi}^{-1}\mathbf{h}_{\varphi\psi}
\]


Permuting back: \[\bar{\mathbf{C}}^{-1} = \begin{pmatrix}
\frac{1}{EA} + \CentroidLocation^\top(E\IntRGoRG)^{-1}\CentroidLocation & \mathbf{0}^\top & 0 & -\CentroidLocation^\top(E\IntRGoRG)^{-1} \\
\mathbf{0} & \frac{1}{G}(\mathbf{H}_{\psi\psi}^{-1} + \frac{1}{\tilde{\Jsv}}\bm{\eta}\otimes\bm{\eta}) & -\frac{1}{G\tilde{\Jsv}}\bm{\eta} & \mathbf{0} \\
0 & -\frac{1}{G\tilde{\Jsv}}\bm{\eta}^\top & \frac{1}{G\tilde{\Jsv}} & \mathbf{0}^\top \\
-(E\IntRGoRG)^{-1}\CentroidLocation & \mathbf{0} & \mathbf{0} & (E\IntRGoRG)^{-1}
\end{pmatrix}\]

\textbf{Final result:}

\[
\boxed{\mathbf{X}^{-1} = \begin{pmatrix}
1 & \mathbf{0}^\top & 0 & \mathbf{0}^\top \\[4pt]
\mathbf{0} & \frac{E}{G}(\mathbf{H}_{\psi\psi}^{-1} + \frac{1}{\tilde{\Jsv}}\bm{\eta}\otimes\bm{\eta})\IntRGoRG & -\Jsv\mathbf{H}_{\psi\psi}^{-1}\CenterShearLocation - \frac{\Jsv}{\tilde{\Jsv}}\bm{\eta}(1 + \bm{\eta}^\top\CenterShearLocation) & \mathbf{0} \\[4pt]
0 & -\frac{E}{G\tilde{\Jsv}}\bm{\eta}^\top\IntRGoRG & \frac{\Jsv}{\tilde{\Jsv}}(1 + \bm{\eta}^\top\CenterShearLocation) & \mathbf{0}^\top \\[4pt]
\mathbf{0} & \mathbf{0} & \mathbf{0} & -\FrameBasis^\times
\end{pmatrix}}\]

where 
\(\bm{\eta}\) is
the shear-torsion coupling vector \eqref{eq:def-eta}, 
\(\tilde{\Jsv}\)
is the reduced torsion constant \eqref{eq}, and 
\(\mathbf{h}_{\varphi\psi}\) from \eqref{eq}.

\subsubsection{Structure of X}

\textbf{Local implementability:} Yes, automatically!

By construction, we defined \(\bm{e} = \mathbf{X}\hat{\bm{e}}\) where
\(\hat{\bm{e}}(x) = \mathbf{L}(x)\bm{p} = (\mathbf{L}_0 + x\mathbf{L}_1)\bm{p}\).
Therefore:
\[\TraceDeDpRoot = \mathbf{X}\mathbf{L}_0, \qquad \TraceDeDpOne = \mathbf{X}\mathbf{L}_1\]

The implementability condition
\(\TraceDeDpOne = \TraceDeDpRoot\mathbf{L}_0^{-1}\mathbf{L}_1\) is
satisfied identically:
\[\TraceDeDpRoot\mathbf{L}_0^{-1}\mathbf{L}_1 = \mathbf{X}\mathbf{L}_0\mathbf{L}_0^{-1}\mathbf{L}_1 = \mathbf{X}\mathbf{L}_1 = \TraceDeDpOne\]


\textbf{Plane-section property:} No!

To check this, we need \(\TraceDeDpRoot = \mathbf{X}\mathbf{L}_0\).
Looking at the \((5\text{-}6, 5\text{-}6)\) block of
\(\mathbf{X}^{-1}\):
\[[\mathbf{X}^{-1}]_{(5:6, 5:6)} = -\FrameBasis^\times\]

Inverting:
\([\mathbf{X}]_{(5:6, 5:6)} = (-\FrameBasis^\times)^{-1} = \FrameBasis^\times\)

Since \([\mathbf{L}_0]_{(5:6, 2:3)} = \oneS\) (the
\(\bm{a}_\Omega \mapsto \bm{\epsilon}_1\) block), we get:
\[[\TraceDeDpRoot]_{(5:6, 2:3)} = \FrameBasis^\times \cdot \oneS = \FrameBasis^\times\]

But the plane-section property requires
\([\TraceDeDpRoot]_{(5:6, 2:3)} = -\FrameBasis^\times\) (curvature from
bending should be
\(\bm{\kappa}_\perp = -\FrameBasis \times \bm{a}_\Omega\)).

\textbf{The sign is wrong.} The energy-conjugate trace does not satisfy
the plane-section property.

\begin{center}\rule{0.5\linewidth}{0.5pt}\end{center}


You're absolutely right to push on this. Let me think carefully about
what we actually have.

\textbf{The core issue:}

We derived \(\mathbf{X}^{-1}\) from the energy conjugacy condition at a
single section. But the resultants themselves depend on \(x\)!
Specifically, for a Saint-Venant solution:

\begin{itemize}
\tightlist
\item
  \(N, \bm{F}, T\) are constant along the beam
\item
  But \(\bm{M}(x) = \bm{M}_0 + x(\FrameBasis \times \bm{F})\) varies
  linearly
\end{itemize}

This is just equilibrium: \(\bm{M}' = \FrameBasis \times \bm{F}\).

\textbf{Let me redo the energy conjugacy condition properly:}

At section \(x\), energy conjugacy requires:
\[\bm{s}(x)^\top \delta\bm{e}(x) = \int_\Section \bm{\sigma}^\top \delta\bm{\varepsilon} \dA = \delta\hat{\bm{e}}(x)^\top \bar{\mathbf{C}} \hat{\bm{e}}(x)\]

If we assume \(\bm{e}(x) = \mathbf{X}\hat{\bm{e}}(x)\) with
\(\mathbf{X}\) independent of \(x\), then:
\[\mathbf{X}^\top \bm{s}(x) = \bar{\mathbf{C}} \hat{\bm{e}}(x)\]

Now express both sides in terms of \(\bm{p}\): -
\(\hat{\bm{e}}(x) = (\mathbf{L}_0 + x\mathbf{L}_1)\bm{p}\) -
\(\bm{s}(x) = -(\IesanStiff + x\mathbf{R}_1)\bm{p}\)

where \(\mathbf{R}_1\) captures the moment variation:
\[\mathbf{R}_1 = \begin{pmatrix}
0 & \mathbf{0}^\top & 0 & \mathbf{0}^\top \\
\mathbf{0} & \mathbf{0} & \mathbf{0} & \mathbf{0} \\
0 & \mathbf{0}^\top & 0 & \mathbf{0}^\top \\
\mathbf{0} & \mathbf{0} & \mathbf{0} & -\FrameBasis^\times E\IntRGoRG
\end{pmatrix}\]

The energy conjugacy condition becomes:
\[-\mathbf{X}^\top (\IesanStiff + x\mathbf{R}_1)\bm{p} = \bar{\mathbf{C}} (\mathbf{L}_0 + x\mathbf{L}_1)\bm{p}\]

Matching powers of \(x\): - \(x^0\):
\(-\mathbf{X}^\top \IesanStiff = \bar{\mathbf{C}} \mathbf{L}_0\) -
\(x^1\):
\(-\mathbf{X}^\top \mathbf{R}_1 = \bar{\mathbf{C}} \mathbf{L}_1\)

From \(x^0\):
\(\mathbf{X}^\top = -\bar{\mathbf{C}} \mathbf{L}_0 \IesanStiff^{-1}\)

Substituting into \(x^1\):
\[\bar{\mathbf{C}} \mathbf{L}_0 \IesanStiff^{-1} \mathbf{R}_1 = \bar{\mathbf{C}} \mathbf{L}_1\]

\textbf{This is a consistency condition.} If it holds, the
energy-conjugate trace with constant \(\mathbf{X}\) is both
implementable and exact. If not, we cannot have both.

\subsubsection{Consistency condition:}

We need: \(\mathbf{L}_0 \IesanStiff^{-1} \mathbf{R}_1 = \mathbf{L}_1\)
or equivalently:
\(\IesanStiff^{-1} \mathbf{R}_1 = \mathbf{L}_0^{-1} \mathbf{L}_1\)

We already computed: \[\mathbf{L}_0^{-1}\mathbf{L}_1 = \begin{pmatrix}
0 & \mathbf{0}^\top & 0 & -\CentroidLocation^\top \\
\mathbf{0} & \mathbf{0} & \mathbf{0} & \mathbf{0} \\
0 & \mathbf{0}^\top & 0 & \mathbf{0}^\top \\
\mathbf{0} & \mathbf{0} & \mathbf{0} & \oneS
\end{pmatrix}\]

Now I need \(\IesanStiff^{-1}\mathbf{R}_1\). Recall:
\[
\IesanStiff = \begin{pmatrix}
EA & EA\CentroidLocation^\top & 0 & \mathbf{0}^\top \\
\mathbf{0} & \mathbf{0} & \mathbf{0} & E\IntRGoRG \\
0 & \mathbf{0}^\top & G\Jsv & \mathbf{0}^\top \\
EA(\FrameBasis\times\CentroidLocation) & E\FrameBasis^\times\IntRoR & \mathbf{0} & \mathbf{0}
\end{pmatrix},
\qquad 
\mathbf{R}_1 = \begin{pmatrix}
0 & \mathbf{0}^\top & 0 & \mathbf{0}^\top \\
\mathbf{0} & \mathbf{0} & \mathbf{0} & \mathbf{0} \\
0 & \mathbf{0}^\top & 0 & \mathbf{0}^\top \\
\mathbf{0} & \mathbf{0} & \mathbf{0} & -\FrameBasis^\times E\IntRGoRG
\end{pmatrix}
\]


The relevant part of \(\IesanStiff^{-1}\) is
\([\IesanStiff^{-1}]_{(2,4)} = (E\IntRGoRG)^{-1}\).

So:
\[[\IesanStiff^{-1}\mathbf{R}_1]_{(2,4)} = (E\IntRGoRG)^{-1} \cdot (-\FrameBasis^\times E\IntRGoRG) = -\FrameBasis^\times\]

But \([\mathbf{L}_0^{-1}\mathbf{L}_1]_{(2,4)} = \mathbf{0}\).

\textbf{The consistency condition fails!}


The energy-conjugate trace as derived gives
\(\TraceDeDpOne \neq \TraceDeDpRoot \mathbf{L}_0^{-1}\mathbf{L}_1\), so it
has nonzero locality defect. Or if we force implementability by setting
\(\TraceDeDpOne = \TraceDeDpRoot \mathbf{L}_0^{-1}\mathbf{L}_1\), then it's
not exactly energy-conjugate at all \(x\).

Let me compute these systematically.

\subsubsection{Structure of X:}

From the block structure of \(\mathbf{X}^{-1}\):
\[\mathbf{X}^{-1} = \begin{pmatrix}
1 & \mathbf{0}^\top & 0 & \mathbf{0}^\top \\
\mathbf{0} & \mathbf{A} & \mathbf{b} & \mathbf{0} \\
0 & \mathbf{c}^\top & d & \mathbf{0}^\top \\
\mathbf{0} & \mathbf{0} & \mathbf{0} & -\FrameBasis^\times
\end{pmatrix},
\qquad\text{and}\qquad 
\mathbf{X} = \begin{pmatrix}
1 & \mathbf{0}^\top & 0 & \mathbf{0}^\top \\
\mathbf{0} & \tilde{\mathbf{A}} & \tilde{\mathbf{b}} & \mathbf{0} \\
0 & \tilde{\mathbf{c}}^\top & \tilde{d} & \mathbf{0}^\top \\
\mathbf{0} & \mathbf{0} & \mathbf{0} & \FrameBasis^\times
\end{pmatrix}
\]

where: 
\begin{itemize}
\item 
\(\mathbf{A} := \frac{E}{G}\Big(\mathbf{H}_{\psi\psi}^{-1} + \frac{\bm{\eta}\otimes\bm{\eta}}{\tilde{\Jsv}}\Big)\IntRGoRG\)
\item
\(\mathbf{b} := -\Jsv\mathbf{H}_{\psi\psi}^{-1}\CenterShearLocation - \frac{\Jsv}{\tilde{\Jsv}}\bm{\eta}(1 + \bm{\eta}^\top\CenterShearLocation)\)
\item \(\mathbf{c}^\top := \frac{E}{G\tilde{\Jsv}}\bm{\eta}^\top\IntRGoRG\) -
\(d := -\frac{\Jsv}{\tilde{\Jsv}}(1 + \bm{\eta}^\top\CenterShearLocation)\)
\end{itemize}


and
\(\begin{pmatrix} \tilde{\mathbf{A}} & \tilde{\mathbf{b}} \\ \tilde{\mathbf{c}}^\top & \tilde{d} \end{pmatrix} = \begin{pmatrix} \mathbf{A} & \mathbf{b} \\ \mathbf{c}^\top & d \end{pmatrix}^{-1}\).

Let me invert the central \(3 \times 3\) block systematically.

\textbf{The central block of \(\mathbf{X}^{-1}\):}
\[\begin{pmatrix} \mathbf{A} & \mathbf{b} \\ \mathbf{c}^\top & d \end{pmatrix}\]

with: -
\(\mathbf{A} = \frac{E}{G}\Big(\mathbf{H}_{\psi\psi}^{-1} + \frac{\bm{\eta}\bm{\eta}^\top}{\tilde{\Jsv}}\Big)\IntRGoRG\)
-
\(\mathbf{b} = -\Jsv\mathbf{H}_{\psi\psi}^{-1}\CenterShearLocation - \frac{\Jsv}{\tilde{\Jsv}}\bm{\eta}(1 + \bm{\eta}^\top\CenterShearLocation)\)
- \(\mathbf{c}^\top = \frac{E}{G\tilde{\Jsv}}\bm{\eta}^\top\IntRGoRG\) -
\(d = -\frac{\Jsv}{\tilde{\Jsv}}(1 + \bm{\eta}^\top\CenterShearLocation)\)

\textbf{Step 1: Compute \(\mathbf{A}^{-1}\).}

Using Sherman-Morrison on
\(\mathbf{H}_{\psi\psi}^{-1} + \frac{\bm{\eta}\bm{\eta}^\top}{\tilde{\Jsv}}\):

With
\(\bm{\eta}^\top\mathbf{H}_{\psi\psi}\bm{\eta} = \mathbf{h}_{\varphi\psi}^\top\mathbf{H}_{\psi\psi}^{-1}\mathbf{h}_{\varphi\psi} = \Jsv - \tilde{\Jsv}\):

\[\Big(\mathbf{H}_{\psi\psi}^{-1} + \frac{\bm{\eta}\bm{\eta}^\top}{\tilde{\Jsv}}\Big)^{-1} = \mathbf{H}_{\psi\psi} - \frac{\mathbf{H}_{\psi\psi}\bm{\eta}\bm{\eta}^\top\mathbf{H}_{\psi\psi}}{\tilde{\Jsv} + \Jsv - \tilde{\Jsv}} = \mathbf{H}_{\psi\psi} - \frac{\mathbf{h}_{\varphi\psi}\mathbf{h}_{\varphi\psi}^\top}{\Jsv}\]

Therefore:
\[\boxed{\mathbf{A}^{-1} = \frac{G}{E}\IntRGoRG^{-1}\Big(\mathbf{H}_{\psi\psi} - \frac{\mathbf{h}_{\varphi\psi}\mathbf{h}_{\varphi\psi}^\top}{\Jsv}\Big)}\]

\textbf{Step 2: Compute \(\mathbf{A}^{-1}\mathbf{b}\).}

Recall
\(\bm{\xi} := \mathbf{h}_{\varphi\psi} + \Jsv\CenterShearLocation = \frac{\nu}{2}\bm{\zeta} - \nu\bm{\lambda}\)
(the Poisson-torsion coupling).

After lengthy but straightforward algebra:
\[\Big(\mathbf{H}_{\psi\psi} - \frac{\mathbf{h}_{\varphi\psi}\mathbf{h}_{\varphi\psi}^\top}{\Jsv}\Big)\Big(\Jsv\mathbf{H}_{\psi\psi}^{-1}\CenterShearLocation + \frac{\Jsv\sigma}{\tilde{\Jsv}}\bm{\eta}\Big) = \Jsv\CenterShearLocation + \mathbf{h}_{\varphi\psi} = \bm{\xi}\]

where \(\sigma := 1 + \bm{\eta}^\top\CenterShearLocation\).

Therefore:
\[\boxed{\mathbf{A}^{-1}\mathbf{b} = -\frac{G}{E}\IntRGoRG^{-1}\bm{\xi}}\]

\textbf{Step 3: Compute \(\mathbf{c}^\top\mathbf{A}^{-1}\).}

Using
\(\bm{\eta}^\top\mathbf{H}_{\psi\psi} = \mathbf{h}_{\varphi\psi}^\top\)
and \(\bm{\eta}^\top\mathbf{h}_{\varphi\psi} = \Jsv - \tilde{\Jsv}\):

\[\mathbf{c}^\top\mathbf{A}^{-1} = \frac{1}{\tilde{\Jsv}}\bm{\eta}^\top\Big(\mathbf{H}_{\psi\psi} - \frac{\mathbf{h}_{\varphi\psi}\mathbf{h}_{\varphi\psi}^\top}{\Jsv}\Big) = \frac{1}{\tilde{\Jsv}}\Big(\mathbf{h}_{\varphi\psi}^\top - \frac{\Jsv - \tilde{\Jsv}}{\Jsv}\mathbf{h}_{\varphi\psi}^\top\Big)\]

Therefore:
\[\boxed{\mathbf{c}^\top\mathbf{A}^{-1} = \frac{\mathbf{h}_{\varphi\psi}^\top}{\Jsv}}\]

\textbf{Step 4: Compute the Schur complement
\(S = d - \mathbf{c}^\top\mathbf{A}^{-1}\mathbf{b}\).}

Define
\(\xi_{\Jsv} := \mathbf{h}_{\varphi\psi}^\top\IntRGoRG^{-1}\bm{\xi}\)  
(a
scalar coupling Poisson-torsion to bending stiffness).

\[\mathbf{c}^\top\mathbf{A}^{-1}\mathbf{b} = -\frac{G}{E\Jsv}\xi_{\Jsv}\]

\[S = -\frac{\Jsv\sigma}{\tilde{\Jsv}} + \frac{G\xi_{\Jsv}}{E\Jsv} = -\frac{E\Jsv^2\sigma - G\tilde{\Jsv}\xi_{\Jsv}}{E\Jsv\tilde{\Jsv}}\]

Define: \[\boxed{\Delta := E\Jsv^2\sigma - G\tilde{\Jsv}\xi_{\Jsv}}\]

Then \(S = -\frac{\Delta}{E\Jsv\tilde{\Jsv}}\) and:

\[
\begin{aligned}
\tilde{d} = S^{-1} &= -\frac{E\Jsv\tilde{\Jsv}}{\Delta}\\
\tilde{\mathbf{b}} &= -\mathbf{A}^{-1}\mathbf{b}\,\tilde{d} = -\frac{G\Jsv\tilde{\Jsv}}{\Delta}\IntRGoRG^{-1}\bm{\xi} \\
\tilde{\mathbf{c}}^\top &= -\tilde{d}\,\mathbf{c}^\top\mathbf{A}^{-1} = \frac{E\tilde{\Jsv}}{\Delta}\mathbf{h}_{\varphi\psi}^\top \\
\tilde{\mathbf{A}} &= \mathbf{A}^{-1} - \frac{\tilde{\mathbf{b}}\tilde{\mathbf{c}}^\top}{\tilde{d}} \\
&= \frac{G}{E}\IntRGoRG^{-1}\Big(\mathbf{H}_{\psi\psi} - \frac{\mathbf{h}_{\varphi\psi}\mathbf{h}_{\varphi\psi}^\top}{\Jsv} + \frac{G\tilde{\Jsv}}{\Delta}\bm{\xi}\mathbf{h}_{\varphi\psi}^\top\Big)
\end{aligned}
\]

\subsection{Trace Matrices}

\textbf{Computing \(\TraceDeDpRoot = \mathbf{X}\mathbf{L}_0\):}

With: \[\mathbf{L}_0 = \begin{pmatrix}
1 & \mathbf{0}^\top & 0 & \mathbf{0}^\top \\
\mathbf{0} & \mathbf{0} & \mathbf{0} & \oneS \\
0 & \mathbf{0}^\top & 1 & \CenterShearLocation^\top \\
\mathbf{0} & \oneS & \mathbf{0} & \mathbf{0}
\end{pmatrix},
\qquad
\mathbf{L}_1 = \begin{pmatrix}
0 & \mathbf{0}^\top & 0 & -\CentroidLocation^\top \\
\mathbf{0} & \mathbf{0} & \mathbf{0} & \mathbf{0} \\
0 & \mathbf{0}^\top & 0 & \mathbf{0}^\top \\
\mathbf{0} & \mathbf{0} & \mathbf{0} & \oneS
\end{pmatrix}
\]

Therefore: \[
\boxed{
\TraceDeDpRoot = \begin{pmatrix}
1 & \mathbf{0}^\top & 0 & \mathbf{0}^\top \\[4pt]
\mathbf{0} & \mathbf{0} & \tilde{\mathbf{b}} & \tilde{\mathbf{A}} + \tilde{\mathbf{b}}\CenterShearLocation^\top \\[4pt]
0 & \mathbf{0}^\top & \tilde{d} & \tilde{\mathbf{c}}^\top + \tilde{d}\CenterShearLocation^\top \\[4pt]
\mathbf{0} & \FrameBasis^\times & \mathbf{0} & \mathbf{0}
\end{pmatrix},
\qquad
\TraceDeDpOne = \begin{pmatrix}
0 & \mathbf{0}^\top & 0 & -\CentroidLocation^\top \\[4pt]
\mathbf{0} & \mathbf{0} & \mathbf{0} & \mathbf{0} \\[4pt]
0 & \mathbf{0}^\top & 0 & \mathbf{0}^\top \\[4pt]
\mathbf{0} & \mathbf{0} & \mathbf{0} & \FrameBasis^\times
\end{pmatrix}}
\]


\textbf{Physical interpretation:}

The fact that
it corresponds to a trace center at the origin is a \emph{consequence}
of the variational definition, not an input.

Specifically, the axial strain \(\epsilon\) is energy-conjugate
to the axial force \(N\), not to any moment. Pure bending
(\(\bm{a}_\Omega \neq 0\) with \(a_\varepsilon = 0\)) produces zero
axial force, so energy conjugacy requires
\(\delta\epsilon^{(a)} \cdot N = 0\), which is satisfied if
\(\epsilon^{(a)} = 0\).

\subsection{Trace Matrices (Corrected)}

Looking back at \([\mathbf{X}^{-1}]_{(5:6, 5:6)} = \FrameBasis^\times\),
we have:
\[[\mathbf{X}]_{(5:6, 5:6)} = (\FrameBasis^\times)^{-1} = -\FrameBasis^\times
\]


The corrected trace matrices are:

\[
\boxed{\TraceDeDpRoot = \begin{pmatrix}
1 & \mathbf{0}^\top & 0 & \mathbf{0}^\top \\[4pt]
\mathbf{0} & \mathbf{0} & \tilde{\mathbf{b}} & \tilde{\mathbf{A}} + \tilde{\mathbf{b}}\CenterShearLocation^\top \\[4pt]
0 & \mathbf{0}^\top & \tilde{d} & \tilde{\mathbf{c}}^\top + \tilde{d}\CenterShearLocation^\top \\[4pt]
\mathbf{0} & -\FrameBasis^\times & \mathbf{0} & \mathbf{0}
\end{pmatrix}
\qquad 
\TraceDeDpOne = \begin{pmatrix}
0 & \mathbf{0}^\top & 0 & -\CentroidLocation^\top \\[4pt]
\mathbf{0} & \mathbf{0} & \mathbf{0} & \mathbf{0} \\[4pt]
0 & \mathbf{0}^\top & 0 & \mathbf{0}^\top \\[4pt]
\mathbf{0} & \mathbf{0} & \mathbf{0} & -\FrameBasis^\times
\end{pmatrix}}\]

Now the \((4, 2:3)\) block of \(\TraceDeDpRoot\) is
\(-\FrameBasis^\times\), consistent with the plane-section property
\(\bm{\kappa}_\perp = -\FrameBasis \times \bm{a}_\Omega\).

\subsection{Correspondance}

Let me reconcile the two derivations carefully.

\subsubsection{The shear-flexure
relation}\label{the-shear-flexure-relation}

Your earlier derivation gives:
\[\bm{\gamma}^{\rm sh} = \BishearTensorRB^{-\top}\mathbf{H}\,\bm{b}_{\Section} \quad \Longleftrightarrow \quad \bm{b}_{\Section} = \mathbf{H}^{-1}\BishearTensorRB^\top\,\bm{\gamma}^{\rm sh}\]

From my \(\mathbf{X}^{-1}\), the \((2,2)\) block (ignoring torsion
coupling) gives:
\[\bm{\beta} = \frac{E}{G}\mathbf{H}_{\psi\psi}^{-1}\IntRGoRG\,\bm{\gamma}_\perp\]

\textbf{Reconciling via equilibrium:}

The key insight is that \(\BishearTensorRB\) (defined from the shear
stress distribution) equals \(E\IntRGoRG\) by equilibrium. Under pure
flexure loading \(\bm{b}_{\Section}\), the shear force is:
\[\bm{F} = \int_\Section G\,\bm{\varepsilon}_\gamma^{(b)} \dA = G\int_\Section \bm{\tau}_\psi \dA \cdot \bm{b}_{\Section} = \BishearTensorRB\,\bm{b}_{\Section}\]

But from the Ieşan equilibrium relation:
\[\bm{F} = E\IntRGoRG\,\bm{b}_{\Section}\]

Therefore: \[\boxed{\BishearTensorRB = E\IntRGoRG}\]

\textbf{Verifying consistency:}

Substituting into your formula:
\[\bm{b}_{\Section} = \mathbf{H}^{-1}\BishearTensorRB^\top\,\bm{\gamma}^{\rm sh} = (G\mathbf{H}_{\psi\psi})^{-1}(E\IntRGoRG)^\top\,\bm{\gamma}^{\rm sh} = \frac{E}{G}\mathbf{H}_{\psi\psi}^{-1}\IntRGoRG\,\bm{\gamma}^{\rm sh}\]

This matches my \([\mathbf{X}^{-1}]_{(2,2)}\) block exactly (when
torsion coupling vanishes). ✓

\textbf{The effective shear stiffness:}

Your formula:
\[\mathbf{K}_{\rm eff} = \BishearTensorRB\,\mathbf{H}^{-1}\BishearTensorRB^\top = (E\IntRGoRG)\,(G\mathbf{H}_{\psi\psi})^{-1}\,(E\IntRGoRG)^\top = \frac{E^2}{G}\IntRGoRG\,\mathbf{H}_{\psi\psi}^{-1}\,\IntRGoRG\]

And the shear force-strain relation:
\[\bm{F} = \mathbf{K}_{\rm eff}\,\bm{\gamma}^{\rm sh}\]

\textbf{Including torsion coupling:}

When \(\CenterShearLocation \neq \mathbf{0}\) or \(\bm{\eta} \neq \mathbf{0}\),
the full \((2,2)\) block of \(\mathbf{X}^{-1}\) is:
\[[\mathbf{X}^{-1}]_{(2,2)} = \frac{E}{G}\Big(\mathbf{H}_{\psi\psi}^{-1} + \frac{\bm{\eta}\bm{\eta}^\top}{\tilde{\Jsv}}\Big)\IntRGoRG\]


\subsection{Kinematic fields}

By integrating the strains.

\textbf{From curvature to rotation:}

We have \(\bm{\kappa} = \bm{\theta}'\), so:
\[\vartheta(x) = \vartheta(0) + \int_0^x \varkappa(s)\,ds\]
\[\bm{\theta}_\perp(x) = \bm{\theta}_\perp(0) + \int_0^x \bm{\kappa}_\perp(s)\,ds\]

From the trace matrices:
\[\varkappa = \tilde{d}\,a_\vartheta + (\tilde{\mathbf{c}}^\top + \tilde{d}\CenterShearLocation^\top)\bm{b}_{\Section} \quad \text{(constant in } x\text{)}\]
\[\bm{\kappa}_\perp(x) = -\FrameBasis^\times\bm{a}_\Omega - x\FrameBasis^\times\bm{b}_{\Section}\]

Integrating:
\[\vartheta(x) = \vartheta(0) + x\big(\tilde{d}\,a_\vartheta + (\tilde{\mathbf{c}}^\top + \tilde{d}\CenterShearLocation^\top)\bm{b}_{\Section}\big)\]
\[\bm{\theta}_\perp(x) = \bm{\theta}_\perp(0) - x\FrameBasis^\times\bm{a}_\Omega - \frac{1}{2}x^2\FrameBasis^\times\bm{b}_{\Section}\]

\textbf{From shear strain to displacement:}

We have \(\bm{\gamma} = \bm{u}' + \FrameBasis \times \bm{\theta}\).

For the axial part, \((\FrameBasis \times \bm{\theta})_\parallel = 0\),
so:
\[u_\parallel(x) = u_\parallel(0) + \int_0^x \epsilon(s)\,ds = u_\parallel(0) + a_\varepsilon x - \frac{1}{2}x^2\CentroidLocation^\top\bm{b}_{\Section}\]

For the transverse part,
\((\FrameBasis \times \bm{\theta})_\perp = \FrameBasis^\times\bm{\theta}_\perp\),
so:
\[\bm{u}_\perp' = \bm{\gamma}_\perp - \FrameBasis^\times\bm{\theta}_\perp\]

With
\(\bm{\gamma}_\perp = \tilde{\mathbf{b}}\,a_\vartheta + (\tilde{\mathbf{A}} + \tilde{\mathbf{b}}\CenterShearLocation^\top)\bm{b}_{\Section}\)
(constant):
\[\bm{u}_\perp(x) = \bm{u}_\perp(0) + x\bm{\gamma}_\perp - \FrameBasis^\times\int_0^x \bm{\theta}_\perp(s)\,ds\]

Computing the integral:
\[\int_0^x \bm{\theta}_\perp(s)\,ds = x\bm{\theta}_\perp(0) - \frac{1}{2}x^2\FrameBasis^\times\bm{a}_\Omega - \frac{1}{6}x^3\FrameBasis^\times\bm{b}_{\Section}\]

Using \((\FrameBasis^\times)^2 = -\mathbf{I}_2\):
\[\bm{u}_\perp(x) = \bm{u}_\perp(0) + x\bm{\gamma}_\perp - x\FrameBasis^\times\bm{\theta}_\perp(0) - \frac{1}{2}x^2\bm{a}_\Omega - \frac{1}{6}x^3\bm{b}_{\Section}\]

\textbf{Matching to the Ieşan solution:}

The section-rigid part of the Ieşan displacement at \(\bm{r} = \mathbf{0}\)
is:
\[\hat{\bm{u}}_{\rm rigid}(x, \mathbf{0}) = xa_\varepsilon\FrameBasis - \frac{1}{2}x^2(\CentroidLocation \cdot \bm{b}_{\Section})\FrameBasis - \frac{1}{2}x^2\bm{a}_\Omega - \frac{1}{6}x^3\bm{b}_{\Section}\]

The warping contribution at \(\bm{r} = \mathbf{0}\) is:
\[\hat{\bm{u}}_{\rm warp}(x, \mathbf{0}) = \WarpTwistIesan(\mathbf{0})(a_\vartheta + c_\vartheta)\FrameBasis + \bm{W}_{(a)}(\mathbf{0})\bm{a}_\Omega + \big(\bm{\Pi}_{(b)}(\mathbf{0}) + x\bm{\Pi}_{(b)}(\mathbf{0}) + \bm{\WarpShearIesan}(\mathbf{0})\big)\bm{b}_{\Section}\]

Wait---let me be more careful. The warping contributions depend on
whether the warping functions are normalized to vanish at the origin.

\textbf{Key observation:}

For the plane-section property with trace center at origin, we have:
\[\bm{u}(x) = \hat{\bm{u}}(x, \mathbf{0}) - \bm{\Phi}(\mathbf{0})\bm{\alpha}(x)\]

At \(x = 0\):
\[\bm{u}(0) = \hat{\bm{u}}(0, \mathbf{0}) - \bm{\Phi}(\mathbf{0})\bm{\alpha}(0)\]

The Ieşan displacement at \((0, \mathbf{0})\) includes: - Axial:
\(\WarpTwistIesan(\mathbf{0})\tau\,\FrameBasis\) (from torsion warping) -
Transverse:
\(\bm{W}_{(a)}(\mathbf{0})\bm{a}_\Omega + \bm{\Pi}_{(b)}(\mathbf{0})\bm{b}_{\Section} + \bm{\WarpShearIesan}(\mathbf{0})\bm{b}_{\Section}\)

\textbf{For a doubly-symmetric section with origin at centroid:}

All warping functions satisfy \(\WarpTwistIesan(\mathbf{0}) = 0\),
\(\bm{W}_{(a)}(\mathbf{0}) = \mathbf{0}\), \(\bm{\Pi}_{(b)}(\mathbf{0}) = \mathbf{0}\),
\(\bm{\WarpShearIesan}(\mathbf{0}) = \mathbf{0}\) by symmetry.

Then: \[\bm{u}(0) = \mathbf{0}, \qquad \bm{\theta}(0) = \mathbf{0}\]

The boundary condition \(\bm{u}(0) = \bm{\theta}(0) = \mathbf{0}\) means
``the origin is fixed and the section has no rotation''---a classical
clamped condition.

\textbf{For a general asymmetric section:}

The warping functions do not vanish at the origin. So:
\[\bm{u}(0) = \WarpTwistIesan(\mathbf{0})\tau\,\FrameBasis + \bm{W}_{(a)}(\mathbf{0})\bm{a}_\Omega + (\bm{\Pi}_{(b)}(\mathbf{0}) + \bm{\WarpShearIesan}(\mathbf{0}))\bm{b}_{\Section}\]

This is nonzero in general!

\subsubsection{Physical
interpretation:}\label{physical-interpretation-1}

The energy-conjugate trace with \(\bm{u}(0) = \bm{\theta}(0) = \mathbf{0}\)
does NOT mean ``the origin is fixed.'' It means:

\[\hat{\bm{u}}(0, \mathbf{0}) = \bm{\Phi}(\mathbf{0})\bm{\alpha}(0)\]

The origin can have nonzero displacement, provided it equals the warping
contribution there. This is an \textbf{implicit constraint coupling the
rigid and warping modes}.

Alternatively, if we want a true ``origin fixed'' boundary condition
\(\hat{\bm{u}}(0, \mathbf{0}) = \mathbf{0}\), then:
\[\bm{u}(0) = -\bm{\Phi}(\mathbf{0})\bm{\alpha}(0) \neq \mathbf{0}\]

\textbf{Summary:}

For asymmetric sections, the energy-conjugate trace boundary condition
\(\bm{u}(0) = \bm{\theta}(0) = \mathbf{0}\) does not correspond to fixing a
physical point. The trace center is at the coordinate origin, but the
warping evaluated there creates an offset between the beam-theoretic
``displacement'' and the actual 3D displacement of that point.

\subsection{Summary}


\textbf{Special case: \(\bm{\xi} = \mathbf{0}\) (doubly-symmetric section or
\(\nu = 0\)).}

When the Poisson-torsion coupling vanishes: -
\(\mathbf{h}_{\varphi\psi} = -\Jsv\CenterShearLocation\) -
\(\xi_{\Jsv} = 0\), so \(\Delta = E\Jsv^2\sigma\) -
\(\tilde{\mathbf{b}} = \mathbf{0}\) -
\(\tilde{d} = -\frac{\tilde{\Jsv}}{\Jsv\sigma}\) -
\(\tilde{\mathbf{c}}^\top = -\frac{\tilde{\Jsv}}{\Jsv\sigma}\CenterShearLocation^\top\)
-
\(\tilde{\mathbf{A}} = \frac{G}{E}\IntRGoRG^{-1}\Big(\mathbf{H}_{\psi\psi} - \Jsv\CenterShearLocation\CenterShearLocation^\top\Big)\)

\textbf{Further special case: \(\CenterShearLocation = \mathbf{0}\) (origin
at shear center).}

Then \(\sigma = 1\), \(\bm{\eta} = \mathbf{0}\), \(\tilde{\Jsv} = \Jsv\), and:
- \(\tilde{\mathbf{b}} = \mathbf{0}\) -
\(\tilde{\mathbf{c}}^\top = \mathbf{0}^\top\) - \(\tilde{d} = -1\) -
\(\tilde{\mathbf{A}} = \frac{G}{E}\IntRGoRG^{-1}\mathbf{H}_{\psi\psi}\)

The shear and torsion decouple completely.
